\chapter{Gazebo Operator Runbook}
\label{ch:operations}

\chapteroverview
  {Provides repeatable Gazebo-only procedures for starting, checking,
   commanding, extending, and stopping the supported runtime profiles.}
  {Proceed in dependency order: verify the overlay, prove the base floor,
   operate robots or rails, add deterministic supervision, add visual state,
   enable authorized task execution, and shut down cleanly. Stop at the first
   unmet readiness condition.}

\section{Scope and Safety Preconditions}

This runbook applies only to the Gazebo simulation described by this repository.
It does not authorize connection to physical rails, PLCs, robots, or grippers;
the complete deployment boundary is defined in \cref{ch:handover}. Follow the
sections in dependency order. Do not begin an optional layer until the preceding
layer's clock, services, topics, and authoritative feedback are healthy.

\begin{safetybox}[Completion requires feedback]
A published command is a request, not proof of completion. Require a newer typed
actual-state sample or a newer verified task effect before the next mutation.
The direct rail commands in this chapter bypass the planner and deterministic
supervisor and therefore require an explicitly controlled, empty-rail procedure.
\end{safetybox}

\section{Runbook Source Map}

The following map identifies the implementation and interface sources behind
each operational group.  Commands later in this chapter are operator recipes;
the listed files remain authoritative when behavior or interfaces change.

\begin{longtable}{L{0.27\textwidth} L{0.65\textwidth}}
\caption{Operator-runbook groups and their authoritative source files}
\label{tab:operator-runbook-source-map}\\
\toprule
\tableheader Runbook group & Source file(s) \\
\midrule
\endfirsthead
\toprule
\tableheader Runbook group & Source file(s) \\
\midrule
\endhead
Terminal setup and preflight &
\file{mfja_3rd_floor_bringup/package.xml};
\file{mfja_3rd_floor_bringup/launch/room_315_floor_common.py}. \\ \rowseparator
Base headless smoke and Room 315 profile &
\file{mfja_3rd_floor_bringup/launch/room_315_only.launch.py};
\file{mfja_3rd_floor_bringup/launch/room_315_floor_common.py};
\file{mfja_robot_control_config/launch/room_315_dual_kinematic_shuttles.launch.py}. \\ \rowseparator
Full-floor profiles &
\file{mfja_3rd_floor_bringup/launch/full_floor.launch.py};
\file{mfja_3rd_floor_bringup/launch/room_315_floor_common.py};
\file{mfja_robot_control_config/launch/multi_robot_sim.launch.py};
\file{mfja_robot_control_config/config/robots.yaml}. \\ \rowseparator
Isolated industrial robot &
\file{mfja_3rd_floor_bringup/launch/single_industrial_robot.launch.py};
\file{mfja_robot_control_config/launch/isolated_industrial_robot.launch.py}. \\ \rowseparator
Robot joint and gripper operation &
\file{mfja_robot_control_config/scripts/robot_joint_command.py};
\file{mfja_robot_control_config/scripts/robot_gripper_command.py};
\file{mfja_robot_control_config/config/gripper_command_defaults.yaml}. \\ \rowseparator
Shuttle creation and lifecycle &
\file{mfja_rail_interfaces/srv/AddShuttle.srv};
\file{mfja_rail_interfaces/msg/ShuttleCommand.msg};
\file{mfja_robot_control_config/scripts/room_315_kinematic_shuttle_node.py};
\file{mfja_robot_control_config/scripts/room_315_multi_shuttle.py}. \\ \rowseparator
Switch operation &
\file{mfja_rail_interfaces/msg/SwitchCommand.msg};
\file{mfja_rail_interfaces/msg/SwitchState.msg};
\file{mfja_robot_control_config/scripts/room_315_rail_devices.py}. \\ \rowseparator
Stopper and sensor operation &
\file{mfja_rail_interfaces/msg/StopperCommand.msg};
\file{mfja_rail_interfaces/msg/StopperState.msg};
\file{mfja_rail_interfaces/msg/SensorFeedback.msg};
\file{mfja_robot_control_config/scripts/room_315_rail_devices.py}. \\ \rowseparator
Camera bridge and deterministic supervisor &
\file{mfja_3rd_floor_bringup/launch/room_315_floor_common.py};
\file{mfja_robot_control_config/launch/room_315_perception_and_safety.launch.py};
\file{mfja_robot_control_config/scripts/room_315_rail_safety_supervisor.py}. \\ \rowseparator
V4 visual runtime &
\file{mfja_robot_control_config/launch/room_315_visual_state_runtime.launch.py};
\file{mfja_robot_control_config/scripts/room_315_visual_state_inference_node.py};
\file{mfja_robot_control_config/scripts/room_315_visual_runtime_v4.py};
\file{mfja_robot_control_config/config/room_315_visual_state/visual_state_runtime.yaml}. \\ \rowseparator
Planning and closed-loop task execution &
\file{mfja_robot_control_config/launch/room_315_task_execution.launch.py};
\file{mfja_robot_control_config/scripts/room_315_task_execution_node.py};
\file{mfja_robot_control_config/scripts/room_315_task_execution.py};
\file{mfja_robot_control_config/config/room_315_task_execution/task_execution_runtime.yaml}. \\ \rowseparator
Controlled shutdown &
\file{mfja_3rd_floor_bringup/launch/room_315_floor_common.py};
\file{mfja_robot_control_config/scripts/room_315_rail_safety_supervisor.py};
\file{mfja_robot_control_config/scripts/room_315_task_execution_node.py}. \\
\bottomrule
\end{longtable}

\section{Terminal Setup and Preflight}

Use the same workspace variables in every terminal. These commands do not build
or change runtime state:

\begin{lstlisting}[language=bash]
export MFJA_WS="${MFJA_WS:-$HOME/mfja_ws}"
export MFJA_REPO="${MFJA_REPO:-$MFJA_WS/src/mfja_3rd_floor_gz}"

cd "$MFJA_WS"
source /opt/ros/jazzy/setup.bash
source "$MFJA_WS/install/setup.bash"

git -C "$MFJA_REPO" status --short --branch
ros2 pkg prefix mfja_3rd_floor_bringup
ros2 pkg prefix mfja_robot_control_config
\end{lstlisting}

Before launching, stop any previous high-level floor with Ctrl-C, confirm the
desired overlay, and decide whether startup may clear the configured obstacle
cache. The high-level default is destructive for that one cache file; the
runbook commands below explicitly preserve it with
\code{room315_clear_visual_obstacle_pose_cache:=false}.

\section{Base Headless Smoke}

Terminal A:

\begin{lstlisting}[language=bash]
ros2 launch mfja_3rd_floor_bringup room_315_only.launch.py \
  robots:=none \
  gui:=false \
  start_paused:=false \
  enable_room315_kinematic_shuttles:=true \
  room315_clear_visual_obstacle_pose_cache:=false \
  room315_right_shuttle_count:=1 \
  room315_left_shuttle_count:=0 \
  room315_shuttles_start_enabled:=false
\end{lstlisting}

After five seconds, Terminal B:

\begin{lstlisting}[language=bash]
ros2 topic echo --once /clock
ros2 service list | grep '^/world/room_315_only/'
ros2 topic echo --once /room_315/rails/right/shuttles/state \
  mfja_rail_interfaces/msg/ShuttleState
ros2 topic echo --once /room_315/rails/right/sensors/feedback \
  mfja_rail_interfaces/msg/SensorFeedback
\end{lstlisting}

Expect \code{create}, \code{remove}, and \code{set_pose} world services plus a
stopped right-shuttle state. \code{gui:=false} omits the Gazebo client; camera
rendering can still require working server-side EGL/graphics support.

\section{Supported Simulation Profiles}

\subsection{Room 315, empty rails}

\begin{lstlisting}[language=bash]
ros2 launch mfja_3rd_floor_bringup room_315_only.launch.py \
  robots:=none gui:=true start_paused:=false \
  room315_clear_visual_obstacle_pose_cache:=false \
  room315_right_shuttle_count:=0 \
  room315_left_shuttle_count:=0
\end{lstlisting}

\subsection{Lightweight full floor}

\begin{lstlisting}[language=bash]
ros2 launch mfja_3rd_floor_bringup full_floor.launch.py \
  robots:=none gui:=true start_paused:=false \
  room315_clear_visual_obstacle_pose_cache:=false \
  room315_right_shuttle_count:=0 \
  room315_left_shuttle_count:=0
\end{lstlisting}

The unmodified full-floor default is paused. This example explicitly starts
running so simulation-time controllers advance.

\subsection{Full floor with selected robots}

\begin{lstlisting}[language=bash]
ros2 launch mfja_3rd_floor_bringup full_floor.launch.py \
  robots:=kuka1,staubli1,yaskawa_hc10_1,yaskawa_hc10dt_1,tiago1 \
  gui:=true start_paused:=false \
  room315_clear_visual_obstacle_pose_cache:=false
\end{lstlisting}

\subsection{Isolated industrial robot}

\begin{lstlisting}[language=bash]
ros2 launch mfja_3rd_floor_bringup single_industrial_robot.launch.py \
  robot:=kuka gui:=true start_paused:=false
\end{lstlisting}

Valid short selectors are \code{kuka}, \code{staubli}, \code{hc10}, and
\code{hc10dt}. Use this profile for fast joint/gripper changes.

\section{Robot Joint and Gripper Operation}

Discover the current installed profiles before commanding:

\begin{lstlisting}[language=bash]
ros2 run mfja_robot_control_config robot_joint_command.py --list
ros2 run mfja_robot_control_config robot_gripper_command.py --list
\end{lstlisting}

Preview first:

\begin{lstlisting}[language=bash]
ros2 run mfja_robot_control_config robot_joint_command.py \
  kuka --unit rad \
  --positions 0.6 -1.0 1.1 0.0 0.6 0.0 \
  --dry-run

ros2 run mfja_robot_control_config robot_gripper_command.py \
  kuka open --dry-run
\end{lstlisting}

Publish the verified command:

\begin{lstlisting}[language=bash]
ros2 run mfja_robot_control_config robot_joint_command.py \
  kuka --unit rad \
  --positions 0.6 -1.0 1.1 0.0 0.6 0.0

ros2 run mfja_robot_control_config robot_gripper_command.py kuka open
ros2 run mfja_robot_control_config robot_gripper_command.py kuka close
ros2 run mfja_robot_control_config robot_gripper_command.py kuka open 75
ros2 run mfja_robot_control_config robot_gripper_command.py \
  hc10 --position 0.012
\end{lstlisting}

The selected robot must already exist and expose the expected subscriber. Check
\rosname{/<robot>/joint_states}, command topic type, and trajectory progress.
Industrial grippers only animate jaws; they do not attach payloads.

\section{Direct Rail Operation}

The following expert examples target the right rail. For the left rail, replace
\code{right} and the entity/identity consistently. Begin with an empty-rail
profile. Direct typed publication bypasses the planner and supervisor.

\subsection{Add and inspect a stopped shuttle}

\begin{lstlisting}[language=bash]
ros2 service call /room_315/rails/right/shuttles/add \
  mfja_rail_interfaces/srv/AddShuttle \
  "{name: 'room315_right_shuttle_1', start_slot: '2', \
    speed: 0.2, start_enabled: false}"

ros2 topic echo --once /room_315/rails/right/shuttles/state \
  mfja_rail_interfaces/msg/ShuttleState
\end{lstlisting}

A successful add reply accepts the logical shuttle; visual spawn is
asynchronous. Verify both state and Gazebo rendering when the visual matters.

\subsection{Start, target, stop, reset, and remove}

The mutation commands below are mutually exclusive examples, not one script to
paste from top to bottom. Choose the applicable ON form, OFF scope, or lifecycle
operation, then require a newer typed actual-state sample before issuing another
mutation.

To start on the retained/default route without an exact target, run only:

\begin{lstlisting}[language=bash]
ros2 topic pub --once /room_315/rails/right/shuttles/command \
  mfja_rail_interfaces/msg/ShuttleCommand \
  "{name: 'room315_right_shuttle_1', command: 'ON', speed: 0.2}"
\end{lstlisting}

Alternatively, to start toward one exact reachable target, run only:

\begin{lstlisting}[language=bash]
ros2 topic pub --once /room_315/rails/right/shuttles/command \
  mfja_rail_interfaces/msg/ShuttleCommand \
  "{name: 'room315_right_shuttle_1', command: 'ON', \
    speed: 0.2, target_slot: '3'}"
\end{lstlisting}

After either ON form, require a newer MOVING or guarded WAITING typed state.
To stop one identity, run:

\begin{lstlisting}[language=bash]
ros2 topic pub --once /room_315/rails/right/shuttles/command \
  mfja_rail_interfaces/msg/ShuttleCommand \
  "{name: 'room315_right_shuttle_1', command: 'OFF'}"
\end{lstlisting}

Use the following as an alternative emergency scope, not immediately after the
single-identity command:

\begin{lstlisting}[language=bash]
ros2 topic pub --once /room_315/rails/right/shuttles/command \
  mfja_rail_interfaces/msg/ShuttleCommand \
  "{name: 'ALL', command: 'OFF'}"
\end{lstlisting}

Require a newer DISABLED state after either OFF scope. To reset a faulted or
stopped identity to its allowed origin, first prove that origin unoccupied,
then run:

\begin{lstlisting}[language=bash]
ros2 topic pub --once /room_315/rails/right/shuttles/command \
  mfja_rail_interfaces/msg/ShuttleCommand \
  "{name: 'room315_right_shuttle_1', command: 'RESET'}"
\end{lstlisting}

Removal is a separate lifecycle operation. After confirming that removal is
intended, run:

\begin{lstlisting}[language=bash]
ros2 topic pub --once /room_315/rails/right/shuttles/command \
  mfja_rail_interfaces/msg/ShuttleCommand \
  "{name: 'room315_right_shuttle_1', command: 'REMOVE'}"
\end{lstlisting}

Nonzero \code{speed} is retained configuration, not proof of motion. A RESET or
OFF shuttle still occupies the rail; verify typed state after every command.

\subsection{Switches}

Only change a coordinated switch pair after proving the rail is empty. Then:

\begin{lstlisting}[language=bash]
ros2 topic pub --once /room_315/rails/right/switches/command \
  mfja_rail_interfaces/msg/SwitchCommand \
  "{switches: [{name: 'A1', state: 'INTERIOR'}, \
               {name: 'A2', state: 'INTERIOR'}]}"
\end{lstlisting}

Alternatively, and only after the same empty-rail proof, request EXTERIOR for
every switch:

\begin{lstlisting}[language=bash]
ros2 topic pub --once /room_315/rails/right/switches/command \
  mfja_rail_interfaces/msg/SwitchCommand \
  "{switches: [{name: 'ALL', state: 'EXTERIOR'}]}"
\end{lstlisting}

After the selected request, inspect the typed actual state:

\begin{lstlisting}[language=bash]
ros2 topic echo --once /room_315/rails/right/switches/state \
  mfja_rail_interfaces/msg/SwitchState
\end{lstlisting}

The pair request and \code{ALL} request are alternatives. A command contains
requested state only: after either example, wait until a newer typed state
reports every requested switch at its delayed actual value. Never issue the
next mutation after a fixed sleep or from GUI animation alone.

\subsection{Stoppers and sensors}

To close stopper A1, run:

\begin{lstlisting}[language=bash]
ros2 topic pub --once /room_315/rails/right/stoppers/command \
  mfja_rail_interfaces/msg/StopperCommand \
  "{stoppers: [{name: 'A1', state: '1'}]}"

ros2 topic echo --once /room_315/rails/right/stoppers/state \
  mfja_rail_interfaces/msg/StopperState
\end{lstlisting}

Only after a newer typed state confirms the close, open the stopper when the
procedure calls for clear passage:

\begin{lstlisting}[language=bash]
ros2 topic pub --once /room_315/rails/right/stoppers/command \
  mfja_rail_interfaces/msg/StopperCommand \
  "{stoppers: [{name: 'A1', state: '0'}]}"

ros2 topic echo --once /room_315/rails/right/stoppers/state \
  mfja_rail_interfaces/msg/StopperState
\end{lstlisting}

Inspect occupancy separately:

\begin{lstlisting}[language=bash]
ros2 topic echo --once /room_315/rails/right/sensors/feedback \
  mfja_rail_interfaces/msg/SensorFeedback
\end{lstlisting}

Close and open are separate mutations. Confirm a newer typed stopper state of
\state{1} before relying on the guard, and confirm \state{0} before relying on
clear passage; sensor feedback does not replace the stopper actual-state check.

\section{Camera Bridge and Deterministic Supervisor}

Start the integrated branch with the floor:

\begin{lstlisting}[language=bash]
ros2 launch mfja_3rd_floor_bringup room_315_only.launch.py \
  robots:=none gui:=true start_paused:=false \
  enable_room315_kinematic_shuttles:=true \
  enable_room315_rail_safety_supervisor:=true \
  enable_room315_rgbd_camera_bridge:=true \
  enable_room315_visual_obstacles:=false \
  room315_clear_visual_obstacle_pose_cache:=false \
  room315_right_shuttle_count:=1 \
  room315_left_shuttle_count:=0 \
  room315_shuttles_start_enabled:=false
\end{lstlisting}

Inspect status and issue the ordinary supervised stop:

\begin{lstlisting}[language=bash]
ros2 topic echo --once /room_315/rail_safety/status std_msgs/msg/String
ros2 topic pub --once /room_315/rail_safety/primitive_command std_msgs/msg/String \
  "{data: '{\"action\":\"stop_all\"}'}"
\end{lstlisting}

This does not require a learned model and does not start visual inference or
task execution. If attaching to a floor that was launched without the
perception-and-safety branch, start
\file{room_315_perception_and_safety.launch.py} directly. Do not do that when the
integrated branch is already active; it would duplicate the supervisor and
camera bridge.

\section{V4 Visual Runtime}

\begin{warningbox}[External artifact prerequisite]
Do not run the checked-in visual YAML blindly: it names a host-specific external
bundle. Obtain a complete independently approved promotion bundle, its JSON
promotion manifest, the independently recorded manifest SHA-256, and the
intended visual virtual environment. Never compute a new ``approved'' hash from
an untrusted file merely to pass the check or edit a promoted bundle in place.
\end{warningbox}

With the integrated RGB-D camera bridge already running, the safe shadow template
is:

\begin{lstlisting}[language=bash]
export ROOM315_SHADOW_V4_BUNDLE="$HOME/room315_artifacts/visual_v4_shadow"
ROOM315_SHADOW_V4_MANIFEST="$ROOM315_SHADOW_V4_BUNDLE"
ROOM315_SHADOW_V4_MANIFEST+=/runtime_promotion_manifest.json
export ROOM315_SHADOW_V4_MANIFEST
export ROOM315_SHADOW_V4_MANIFEST_SHA256='<approved-manifest-sha256>'

source "$HOME/.venvs/mfja-visual/bin/activate"
test -f "$ROOM315_SHADOW_V4_MANIFEST"
printf '%s  %s\n' "$ROOM315_SHADOW_V4_MANIFEST_SHA256" \
  "$ROOM315_SHADOW_V4_MANIFEST" | sha256sum --check -

ros2 launch mfja_robot_control_config \
  room_315_visual_state_runtime.launch.py \
  use_sim_time:=true enable_camera_bridge:=false \
  runtime_mode:=shadow \
  v4_promotion_manifest:="$ROOM315_SHADOW_V4_MANIFEST" \
  v4_promotion_manifest_sha256:="$ROOM315_SHADOW_V4_MANIFEST_SHA256" \
  device:=auto dry_run_state_fusion:=true \
  plansys2_update_enabled:=false
\end{lstlisting}

Inspect:

\begin{lstlisting}[language=bash]
ros2 topic list | sort | rg '/room_315/visual_state/shadow_v4'
ros2 topic echo --once \
  /room_315/visual_state/shadow_v4/validation \
  mfja_rail_interfaces/msg/VisualStateObservation
ros2 topic echo \
  /room_315/visual_state/shadow_v4/diagnostics \
  diagnostic_msgs/msg/DiagnosticArray
\end{lstlisting}

Stop the shadow node before starting the independently approved active bundle.
With the integrated camera bridge still running, use this exact first-readiness
template:

\begin{lstlisting}[language=bash]
export ROOM315_ACTIVE_V4_BUNDLE="$HOME/room315_artifacts/visual_v4_active"
ROOM315_ACTIVE_V4_MANIFEST="$ROOM315_ACTIVE_V4_BUNDLE"
ROOM315_ACTIVE_V4_MANIFEST+=/runtime_promotion_manifest.json
export ROOM315_ACTIVE_V4_MANIFEST
export ROOM315_ACTIVE_V4_MANIFEST_SHA256='<approved-manifest-sha256>'

source "$HOME/.venvs/mfja-visual/bin/activate"
test -f "$ROOM315_ACTIVE_V4_MANIFEST"
printf '%s  %s\n' "$ROOM315_ACTIVE_V4_MANIFEST_SHA256" \
  "$ROOM315_ACTIVE_V4_MANIFEST" | sha256sum --check -

ros2 launch mfja_robot_control_config \
  room_315_visual_state_runtime.launch.py \
  use_sim_time:=true enable_camera_bridge:=false \
  runtime_mode:=active \
  v4_promotion_manifest:="$ROOM315_ACTIVE_V4_MANIFEST" \
  v4_promotion_manifest_sha256:="$ROOM315_ACTIVE_V4_MANIFEST_SHA256" \
  device:=auto dry_run_state_fusion:=true \
  plansys2_update_enabled:=false
\end{lstlisting}

In another sourced terminal, require an accepted active observation and inspect
the visual-runtime diagnostic:

\begin{lstlisting}[language=bash]
ros2 topic echo --once \
  /room_315/visual_state/observed_state \
  mfja_rail_interfaces/msg/VisualStateObservation
ros2 topic echo --once /diagnostics \
  diagnostic_msgs/msg/DiagnosticArray
\end{lstlisting}

The observation must have stage \code{fused_observed_state}, current source
stamps, matching schema/checkpoint, \code{accepted}, \code{model_ready},
\code{input_ready}, \code{presence_ready}, and \code{state_fusion_ready} true,
and \code{stale} false. \code{stabilized=false} is valid while the optional
temporal filter is disabled.

Both supervisor and visual launches default to creating RGB bridges. When the
integrated supervisor already owns them, visual runtime must use
\code{enable_camera_bridge:=false}.

\section{Planning and Closed-Loop Task Execution}

Planning process/gateway startup with actuation disabled:

\begin{lstlisting}[language=bash]
ros2 launch mfja_robot_control_config \
  room_315_task_execution.launch.py \
  use_sim_time:=true execution_enabled:=false \
  enable_plansys2:=true external_obstacles_disabled:=true

ros2 lifecycle get /planner
\end{lstlisting}

With \code{execution_enabled=false}, authorization verification is deliberately
bypassed and no authorization or promotion artifact is required. The planner
may be active while the disabled gateway still reports that live visual,
supervisor, or sensor inputs are not ready; that is independent of artifact
availability. For actuation, first create a host-local YAML and edit only its
external paths/hashes:

\begin{lstlisting}[language=bash]
export ROOM315_TASK_EXEC_CONFIG="$HOME/.config/mfja/task_execution_runtime.yaml"
mkdir -p "$(dirname "$ROOM315_TASK_EXEC_CONFIG")"
cp "$MFJA_REPO/mfja_robot_control_config/config/room_315_task_execution/task_execution_runtime.yaml" \
  "$ROOM315_TASK_EXEC_CONFIG"
${EDITOR:-nano} "$ROOM315_TASK_EXEC_CONFIG"
\end{lstlisting}

Set literal absolute paths for \code{task_execution_authorization_path} and
\code{task_execution_promotion_manifest_path}, plus the independently approved
\code{task_execution_authorization_sha256}. Verify that the visual schema and
checkpoint hash match the active observation. Before enabling execution, the
authorization JSON, promotion manifest, manual decision, and every artifact
bound from that manifest must be a regular file, not a symbolic link at its
final path component, and have all user/group/other write bits cleared. The
host-local runtime YAML remains configuration; do not make writable copies of
the immutable artifacts merely to change its paths.

Only after the floor, supervisor, active accepted vision, planner, and artifact
checks are ready:

\begin{lstlisting}[language=bash]
ros2 launch mfja_robot_control_config \
  room_315_task_execution.launch.py \
  use_sim_time:=true \
  runtime_config:="$ROOM315_TASK_EXEC_CONFIG" \
  execution_enabled:=true enable_plansys2:=true \
  external_obstacles_disabled:=true

ros2 lifecycle get /planner
ros2 topic echo /diagnostics diagnostic_msgs/msg/DiagnosticArray
\end{lstlisting}

\code{execution_enabled:=true} causes supervised Gazebo actuation. The current
gateway has no external-obstacle perception provider and supports only the
paired policy
\code{enable_room315_visual_obstacles:=false} with
\code{external_obstacles_disabled:=true}. Setting the gateway value false fails
closed with \code{external_obstacle_perception_not_configured}; task execution
is unavailable while the floor obstacles are enabled.

TaskGoals should normally be published by the confirmatory English CLI. If a
developer publishes JSON directly, first validate it with the same TaskGoal
builder/validator; do not invent fields from memory.

\section{Controlled Shutdown}

\begin{procedurebox}[Shutdown order]
\begin{enumerate}
  \item Stop the active goal or publish supervisor \code{stop_all}.
  \item Require newer DISABLED feedback for every managed shuttle.
  \item Stop the English CLI and separately launched task/visual processes with
        Ctrl-C.
  \item Press Ctrl-C once in the high-level floor terminal and allow children to
        exit; preserve logs if a service/entity error occurred.
  \item Confirm the relevant ROS nodes/Gazebo processes are gone before starting
        another floor. A leftover lock pathname alone is harmless.
\end{enumerate}
\end{procedurebox}

Avoid \code{pkill} in normal operation. The payload batch/review utilities use
broad \code{pkill -f} patterns and restart the ROS daemon; run them only on a
dedicated simulation host after reviewing side effects. Their reset options can
also delete a selected dataset directory. The intent-model setup downloads
about 1.04 GiB and installs packages; keep it inside its isolated environment.

\noindent\textbf{Source files:}
\file{mfja_robot_control_config/scripts/room_315_payload_case_batch_runner.py},
\file{mfja_robot_control_config/scripts/room_315_payload_case_review.py}, and
\file{mfja_robot_control_config/scripts/setup_room315_intent_model.py}.
